
\subsection*{2.2}
\phantomsection
\label{sga3:VIIB:2.2}
\addcontentsline{toc}{subsection}{2.2. Bialgebras over \(\mathbf O_k\)}

Let \(\mathbf H\) be a group of the category of coalgebras over
\(\mathbf O_k\) \textup{(\SGARefLink{sga3:VIIB:1.2.2}{1.2.2})}. For every \(k\)-algebra
\(C\in\operatorname{Alf}/k\), \(\mathbf H(C)\) is therefore endowed with
the structure of a \(C\)-coalgebra in groups \textup{(VII\textsubscript{A}
\SGARefLink{sga3:VIIA:3.2}{3.2}; following Manin, we shall say \emph{bialgebra} instead of
``coalgebra in groups'')}; moreover, if \(\phi:C\to D\) is a morphism of
\(\operatorname{Alf}/k\), the map
\[
  D\otimes_C\mathbf H(C)\longrightarrow\mathbf H(D),
  \qquad
  d\otimes x\longmapsto d\cdot\mathbf H(\phi)(x),
\]
is a homomorphism of bialgebras over \(D\). We shall summarize these
properties by saying that \(\mathbf H\) is a \emph{bialgebra over
\(\mathbf O_k\)}.

It is clear that the functor
\(\mathbf H\mapsto\operatorname{Spf}^*(\mathbf H)\) of \SGARefLink{sga3:VIIB:1.3.5}{1.3.5}
commutes with finite products. It therefore transforms a bialgebra
over \(\mathbf O_k\) into a \(k\)-functor in groups, that is, a functor
from \(\operatorname{Alf}/k\) to the category of groups; and indeed,
for every \(k\)-algebra \(C\) of finite length, the elements \(x\) of
\(\mathbf H(C)\) that belong to \(\operatorname{Spf}^*(\mathbf H)(C)\)
form a group for the multiplication of the algebra \(\mathbf H(C)\)
\textup{(cf.\ VII\textsubscript{A} \SGARefLink{sga3:VIIA:3.2.2}{3.2.2})}: recall that these elements \(x\)
are those such that \(\varepsilon(x)=1\) and \(\Delta(x)=x\otimes x\).
The second condition implies, moreover, the equality
\(\varepsilon(x)=\varepsilon(x)\varepsilon(x)\), hence also
\(\varepsilon(x)=1\) if \(C\) is local and \(x\neq0\).\footnote{Translator's note: the 1970 text has ``hence also
\(\varepsilon(x)=1\) if \(k\) is local''; the reading above is that of the
Gille--Polo re-edition (2011), Exposé~VII\textsubscript{B}. By the counit
property, \(\Delta(x)=x\otimes x\) gives \(x=\varepsilon(x)x\), so the
idempotent \(\varepsilon(x)\) of \(C\) equals \(1\) when \(C\) is local and
\(x\neq0\). That \(k\) be local does not suffice: if \(k\) is a field,
\(\mathbf H=\mathbf O_k\) with its trivial structure \textup{(}\(\mathbf H(C)=C\),
\(\Delta\) and \(\varepsilon\) the canonical isomorphisms\textup{)} and
\(C=k\times k\), the element \(x=(1,0)\) satisfies \(\Delta(x)=x\otimes x\)
and \(\varepsilon(x)=x\neq1\).}
